% Lösungen Einheit 3 von 3 – Fehlerarten und Güte: den Fehler zweiter Art für selbst gewählte Anteile berechnen (p dort wählen, wo die Nullhypothese falsch ist), einordnen und im Sachzusammenhang beschreiben; Mindestanteile für eine Fehlerschranke; die Gütekurve lesen (Ablehnwahrscheinlichkeit in Abhängigkeit von p (zusammenbau v0.1)
\einheitenkopf{Einheit 3 von 3 · Fehlerarten und Güte: den Fehler zweiter Art für selbst gewählte Anteile berechnen (p dort wählen, wo die Nullhypothese falsch ist), einordnen und im Sachzusammenhang beschreiben; Mindestanteile für eine Fehlerschranke; die Gütekurve lesen (Ablehnwahrscheinlichkeit in Abhängigkeit von p}

\erg{13}{a) $p = 0{,}35$ \quad b) $p$ muss kleiner als 0,04 sein: für $p = 0{,}03$ ist $P(X \ge 13) \approx 0{,}7362$, für $p = 0{,}02$ ist $P(X \ge 13) \approx 0{,}2065$; mit dieser Wahrscheinlichkeit wird das Gewinnspiel verlängert, obwohl es sich nicht lohnt \quad c) z. B. $p = 0{,}8$ (dort ist $H_0$ falsch): abgelesen etwa $0{,}79$, Fehler zweiter Art $\approx 1 - 0{,}79 = 0{,}21$ \quad d) $p$ knapp über 0{,}35 wählen, etwa $p = 0{,}36$: $P(X \le 119) \approx 0{,}9160 > 0{,}90$ \quad e) für $p = 0{,}17$ ist $P(X \le 45) \approx 0{,}2003 > 0{,}1$, für $p = 0{,}18$ ist $P(X \le 45) \approx 0{,}0987$, also mindestens 18\,\%; Fehler zweiter Art: es fährt keine zusätzliche Radfähre, obwohl mehr als 12\,\% ein Rad dabeihaben \quad f) An der Grenze $p = 0{,}05$ liest man den Fehler erster Art ab: etwa $0{,}055$. Bei hundert Armbändern wäre $P(Y \le 3) \approx 0{,}2578$; da diese Wahrscheinlichkeit mit wachsendem Umfang fällt, ist der Umfang größer als hundert}
\erg{14}{a) Die Fertigung wird umgestellt, obwohl höchstens der zulässige Anteil ausfällt (unnötige Umstellung); oder sie wird nicht umgestellt, obwohl mehr als der zulässige Anteil ausfällt (Mangel bleibt unentdeckt). Beide beziehen sich auf die Entscheidung über die Fertigung, nicht auf einzelne Lampen.}
\erg{15}{a) Bei $p = 0{,}3$ ist $H_0$ wahr, dort gibt es keinen Fehler zweiter Art; $p$ muss größer als 0,35 sein, etwa $p = 0{,}45$: $P(X \le 19) \approx 0{,}6844$ \quad b) Er ist die Wahrscheinlichkeit des Annahmebereichs unter dem wahren Anteil; je weiter dieser von der Grenze der Nullhypothese entfernt ist, desto seltener fällt das Ergebnis in den Annahmebereich. \quad c) Fehler zweiter Art: $P(X \le 12) \approx 0{,}2612$ – etwa jede vierte solche Anlage bleibt ungewartet; das spricht für eine größere Stichprobe \quad d) 0,6: $\approx 0{,}40$; 0,7: $\approx 0{,}72$; 0,8: $\approx 0{,}94$; 0,9: $\approx 1{,}00$}
